\documentclass[../../../include/open-logic-section]{subfiles}

\begin{document}

\olfileid{sth}{choice}{intro}

\olsection{مقدمه}

در \crefrange{sth:cardinals::chap}{sth:card-arithmetic::chap}، با درنظرگرفتن
کاردینال‌ها به‌عنوان اعداد ترتیبی، نظریه‌ای برای کاردینال‌ها پروراندیم.
آن رویکرد به اصل خوش‌ترتیبی وابسته است. چنان‌که معلوم می‌شود، اصل
خوش‌ترتیبی با اصل دیگری---اصل انتخاب---هم‌ارز است، و دربارهٔ پذیرفتنی‌بودن
آن بحث‌های فلسفی جدی‌ای درگرفته است. پرسش‌های ما در این فصل این‌هاست:
اصل انتخاب چگونه به کار می‌رود، و آیا می‌توان آن را توجیه کرد؟

\end{document}
